\section{Working Record of Intermediate Evidence}
\label{app:intermediate}

This appendix is a working record, not a substitute for the final experimental matrix. The September 12, 2026 report evaluates the common epoch-2 checkpoint for each task, training seed 3072, and 50 development episodes under the revised protocol.

\begin{table}[h]
 \centering
 \begin{tabular}{lrrr}
  \toprule
  Task & P0 & P3 & P4 \\
  \midrule
  Cube & 98 & 98 & 94 \\
  Push-T & 80 & 68 & 78 \\
  Reacher & 28 & 44 & 32 \\
  TwoRoom & 46 & 42 & 56 \\
  \bottomrule
 \end{tabular}
 \caption{Intermediate development success (\%), not final performance. P3/P4 use 64 candidates and 16 flow steps; B verification is split into batches of eight.}
 \label{tab:epoch2}
\end{table}

P4 minus P3 is $-4,+10,-12,+14$ percentage points across the four tasks. The signs are mixed. P4 minus direct P0 execution is $-4,-2,+4,+10$ points. No cross-task improvement, statistical significance, or seed-expansion decision follows from these intermediate values. No latency result is available in the source report.

\section{Training and Evaluation Details to Complete}

\TODO{Appendix artifacts: insert run identity, dataset hashes, exclusion audit, software versions, checkpoint hashes, optimizer schedule, per-seed results, Wilson intervals, paired flips, and timing boundaries.}

The fitting-time predicted-action probability reaches at most 0.5. The inherited non-fitting A/B loss path instead uses predicted actions with probability one. Validation loss is therefore not interchangeable with the training mixture, nor is it used to select an intermediate test-success peak.

\TODO{Figures: architecture with attention and gradient boundaries; candidate coverage versus selected-candidate success; paired failure examples; and accepted success--latency curves. Real execution traces are needed for qualitative outcome figures.}
